\chapter*{General Conclusion and Perspectives}
\addcontentsline{toc}{chapter}{General Conclusion and Perspectives}
\markboth{General Conclusion}{General Conclusion}

The objective of this final year project was to replace a scattered, spreadsheet-based project
steering process with a centralised platform that preserves the richness of the existing financial
model while removing its fragility. That objective is met: PMS covers the whole functional scope
identified during the analysis of the reference workbook, from dynamic authorization to the
earned-value indicator engine.

\section*{Assessment}
The construction proceeded through ten increments, each demonstrable in isolation. Release~1
delivered the technical foundation, authentication and above all the dynamic authorization model.
Release~2 opened operational steering with projects, teams and workload. Release~3 delivered the
core business value: billing, missions, governance and the indicator engine. Release~4
consolidated the interface, verified the security model, and made the platform reproducibly
deployable through containerization and a continuous integration pipeline.

Three design decisions structure the result and are, in retrospect, the ones that mattered most.

The first is the \textbf{dynamic authorization model}. By inserting a permission level between
roles and modules, and by keeping the role--permission matrix in the database, the platform makes
its security policy configurable data rather than compiled code. Granting a new right is an
administrative act, not a release.

The second is the \textbf{hybrid indicator engine}. Rather than choosing between recomputing
indicators on demand and persisting them, the platform does both: live computation guarantees that
no dashboard displays a stale figure, while a frozen monthly snapshot preserves the official state
of each review. This reproduces the strength of the original workbook, a live dashboard with monthly
history, without its fragility.

The third is the \textbf{normalisation of the monthly time series}. Replacing a grid of month
columns with one row per project, resource and month turned a structural limitation of the
spreadsheet into a simple insertion, and made integrity constraints expressible: no duplicate
month, no entry on an unassigned project, no silent rewriting of a validated period.

\section*{Skills Acquired}
Beyond the technologies used, this project required reading a real business model out of an
artefact that was never written as a specification, and arbitrating between what the workbook did,
what the written specifications said, and what the code could reasonably guarantee. It also
required designing for confidentiality: deciding that commercially sensitive amounts should be
computed on read rather than stored is an engineering decision with a business justification, not a
technical preference.

\section*{Limitations and Perspectives}
Several avenues remain open.

\textbf{Automatic time-tracking import.} The workload model was designed to accept a second data
source alongside manual entry, so that effort recorded in an external time-tracking tool could be
imported and reconciled. The mechanism is anticipated in the design but not implemented; it is the
most immediate extension.

\textbf{Reporting and export.} The platform computes and displays the indicators; generating
exportable review documents in the format the company already circulates would shorten the last
manual step of the monthly review.

\textbf{Integration with the information system.} The business rules explicitly anticipate future
integration points with accounting, human resources and enterprise systems. The modular
organisation of the code was chosen with that reversibility in mind: a module can be extracted
behind an interface without restructuring the whole.

\textbf{Scaling the deployment.} The current deployment targets the internal use of one company. A
significant growth in the number of projects or users would justify revisiting the modular monolith,
a decision deliberately taken as reversible rather than permanent.

The platform delivered is therefore both an answer to an immediate operational need and a
foundation on which the company's project steering can continue to evolve.
